\documentclass{situationist-preprint}

\title{Supervaluation and Future Contingents}
\msid{TSM-12-b}
\msfamily{TSM-12}
\msdiscipline{Time, Space and Modality}
\mstype{Companion manuscript}
\mskeywords{future contingents; supervaluation; branching time; settledness; vagueness}

\begin{document}
\maketitle

\begin{abstract}
TSM-12-a argues that future contingents are neither true nor false until they are settled. Supervaluation over the histories through a moment is the natural semantics for this view: it keeps classical logic while denying bivalence. Supervaluation is also used for vagueness, where on the view of MET-18 it measures what can be known about a sharp fact. We argue that the two uses differ in exactly the respect TSM-12-a identifies: in the vague case something grounds a sharp boundary, while relative to a moment nothing grounds the choice of one history. Supertruth at a moment is therefore truth at that moment, not knowable truth, and the gap it leaves is a gap in truth.
\end{abstract}

\section{Supertruth as settled truth}

\citet{Thomason1970} defined a sentence to be true at a moment $m$ if it is true on every history through $m$, false if false on every one, and neither otherwise. This is Definition~2.2 of TSM-12-a: supertruth at $m$ is settled truth at $m$. The semantics has two familiar features. \enquote{There will be a sea battle or there will not} is true at $m$, since each history makes one disjunct true; yet neither disjunct is true at $m$. Excluded middle holds while bivalence fails, and disjunction is not truth-functional at the level of supertruth.

Both features fit the open future. The disjunction is settled, since every history through $m$ makes one disjunct true, while which disjunct holds is not settled. By Theorem~3.2 of TSM-12-a the open disjuncts lack truth value at $m$, and by Proposition~3.3 they are settled at the later moments at which the matter has occurred.

\section{Two uses of supervaluation}

Supervaluation was introduced for vagueness \citep{Fine1975}, and MET-05 uses it to select one cloud among many candidate aggregates. MET-18 holds nonetheless that vague predicates have sharp but unknowable boundaries; on the combination, the admissible precisifications are those not excluded by what can be known about the boundary, and supertruth measures what can be known, while there is a fact about where the boundary lies \citep{Williamson1994}. It may be asked whether the same holds of future contingents: whether supertruth at $m$ measures only what can be known at $m$, while there is a fact about which history will be followed.

It does not, and the reason is Principle~P1 of TSM-12-a. A sharp boundary for \enquote{cloud} or \enquote{heap} is grounded in the use of the word, which is complete when the question is asked; that the boundary is unknowable does not make it ungrounded. A fact about which history through $m$ will be followed would have to be grounded in what is the case relative to $m$: the branching structure, the past of $m$ and the laws. By Principle~P2, nothing in these distinguishes one history from the others \citep{BelnapGreen1994}. There is therefore no fact for supertruth to fall short of.

\begin{proposition}\label{p:gap}
For future contingents, supertruth at $m$ is truth at $m$, and the gap between supertruth and superfalsity is a gap in truth, not in knowledge.
\end{proposition}

\begin{proof}
If a sentence open at $m$ had a truth value at $m$, it would have it in virtue of what is the case relative to $m$ (Principle~P1 of TSM-12-a). Since it is true on some histories through $m$ and false on others, a truth value would require something relative to $m$ that selects one history, which Principle~P2 excludes. Settled sentences are true or false on every history through $m$, and so have the value that supertruth assigns.
\end{proof}

The two uses of supervaluation thus have one form and different interpretations, and the difference is fixed by grounding. Where the fact is grounded but inaccessible, supervaluation describes knowledge; where nothing grounds a fact, it describes truth.

\begin{objection}
The future will be one way. When the battle occurs, it will have been true that it would occur; so there was a fact all along.
\end{objection}

\begin{reply}
At the later moment, \enquote{it was the case that there will be a battle} is settled true, evaluated on the histories through that moment \citep{Thomason1970}. It does not follow that \enquote{there will be a battle} was true at the earlier moment, relative to which several histories remained. The objection evaluates a past-tensed sentence at one moment and draws a conclusion about another. Settledness is moment-relative, as REL-05-b also requires in treating divine foreknowledge.
\end{reply}

\begin{objection}
On assessment-relativism \citep{MacFarlane2003}, a future contingent is true as assessed from a later moment at which it has been settled. That is truth, not mere knowledge.
\end{objection}

\begin{reply}
Truth as assessed later is truth at the later moment of a sentence about an earlier one, which Proposition~3.3 of TSM-12-a already provides. What it does not provide is truth at the moment of utterance, and that is what Theorem~3.2 of TSM-12-a denies.
\end{reply}

\section{Conclusion}

Supervaluation over histories is the semantics of the open future: it keeps classical logic and denies bivalence for future contingents. Unlike supervaluation for vagueness on the view of MET-18, it does not measure the knowability of a fact, because relative to a moment there is no fact about the future to know. Future contingents are neither true nor false until they are settled.

\begin{thebibliography}{99}\small
\bibitem[Belnap and Green(1994)]{BelnapGreen1994} Belnap, N. and Green, M. (1994). Indeterminism and the thin red line. \emph{Philosophical Perspectives}, 8, 365--388.
\bibitem[Fine(1975)]{Fine1975} Fine, K. (1975). Vagueness, truth and logic. \emph{Synthese}, 30(3--4), 265--300.
\bibitem[MacFarlane(2003)]{MacFarlane2003} MacFarlane, J. (2003). Future contingents and relative truth. \emph{The Philosophical Quarterly}, 53(212), 321--336.
\bibitem[Thomason(1970)]{Thomason1970} Thomason, R.~H. (1970). Indeterminist time and truth-value gaps. \emph{Theoria}, 36(3), 264--281.
\bibitem[Williamson(1994)]{Williamson1994} Williamson, T. (1994). \emph{Vagueness}. Routledge.
\end{thebibliography}

\end{document}
